% Research proposal: no new benchmark results or correctness guarantees are claimed.
\documentclass{article}
\usepackage{iclr2027_conference,times}
\usepackage[T1]{fontenc}
\usepackage{amsmath,amssymb,hyperref,url,etoolbox}
\hypersetup{colorlinks=true,linkcolor=blue,citecolor=blue,urlcolor=blue}
\title{CertMem: Evidence Residuals\\for Token-Efficient Agent Memory}
\author{Anonymous authors}
% Label the working draft accurately without changing the official style file.
\makeatletter
\patchcmd{\@maketitle}{Paper under double-blind review}{Manuscript in preparation}{}{}
\makeatother
\begin{document}
\maketitle
\lhead{Working draft: proposed framework; empirical evaluation pending}

\begin{abstract}
Long-term conversational agents must retain detailed histories while limiting the
tokens spent maintaining and consulting memory. Existing systems connect compact
summaries to original evidence, enabling selective recovery of information lost
during compression. However, recoverability alone leaves a coupled allocation
problem: which details should be preserved in compact memory, and which should
be recovered when a question arrives? We propose CertMem, a framework that uses
a shared record of compression omissions to coordinate these decisions. At write
time, CertMem audits candidate summaries against source-grounded checks and
retains a small set of verbatim evidence residuals, prioritizing their predicted
reduction in downstream answering loss per added token. At read time, the same
audit guides selective access to additional source spans, with reassessment after
each expansion. The controller is designed to learn from paired answering
outcomes, while independent validation separates useful diagnostic scores from
claims of answer correctness. We formulate memory efficiency in terms of total
maintenance and query-time token expenditure and specify matched-budget
evaluation on LoCoMo and LongMemEval. The central hypothesis is that preserving
the right omissions before retrieval can reduce repeated source reading and
improve the accuracy--cost frontier beyond summary-based routing alone.
\end{abstract}

\section{Introduction}

A conversational agent may remember that a user discussed a conference in Kyoto
yet forget that the trip was canceled. A later question about the user's travel
plans then requires a distinction that the stored summary no longer expresses.
Retaining the original conversation makes the cancellation recoverable, but
consulting the entire discussion for every related question spends tokens on
information that the summary already preserves. Long-term memory therefore
requires decisions about both what to retain in compact form and how much
original evidence to read. Benchmarks such as LoCoMo and LongMemEval make these
decisions consequential by testing facts, temporal relations, and updates across
extended interaction histories \citep{locomo,longmemeval}.

Recent work provides increasingly effective connections between abstraction and
evidence. RMM stores topic summaries paired with raw dialogue and adapts memory
ranking using citation feedback \citep{rmm}. HiGMem organizes events above
dialogue turns, using event summaries and fact sheets to guide LLM selection of
the turns that support an answer \citep{higmem}. H-MEM routes through domain,
category, memory-trace, and episode layers using indices to subordinate memories
\citep{hmem}. SEEM combines relational facts with episodic event frames and
uses reverse provenance expansion to recover associated source passages
\citep{seem}. Most directly, TierMem answers from summaries when sufficient,
escalates to linked raw pages otherwise, and can consolidate verified findings
for subsequent queries \citep{tiermem}. These advances establish hierarchical
retrieval, provenance links, and selective source access as foundations on which
our work builds.

They also motivate a more specific question: how should information about
\emph{what compression removed} influence both retention and retrieval? A source
pointer identifies where a memory came from; it does not, by itself, describe
which distinctions survived its summary. In the Kyoto example, a brief
cancellation utterance may be more valuable than a longer, fluent description
of the conference. A system can retain that utterance alongside the summary or
recover it during a future query. The two choices have different costs that
depend on how often the information is needed. We hypothesize that explicitly
recording omissions can improve this allocation by letting the writer preserve
inexpensive, consequential details and informing the reader about unresolved
gaps. The hypothesis concerns how compression diagnostics guide both stages.

Evaluating this hypothesis requires distinguishing summary fidelity from task
performance. Preserving many extracted facts can still leave an answer wrong
when a missing condition, attribution, or temporal relation determines the
result. Conversely, a detail omitted from a summary need not matter to the
questions the agent receives. TRUSTMEM already evaluates memory transitions
through coverage, preservation, and faithfulness \citep{trustmem}, while DeMem
formulates memory quality through decision loss induced by compression
\citep{demem}. Our proposed mechanism connects such diagnostic information to a
concrete intervention: selecting short source passages that repair predicted
answering failures and reusing the remaining gaps to guide later evidence access.
Evaluation must include downstream failures that the diagnostic checks did not
anticipate.

We propose \textbf{CertMem}, a framework for coordinating memory compression,
evidence retention, and source expansion under token budgets. For each source
segment, a writer produces candidate summaries at several lengths. A bounded
audit checks source-grounded propositions and selected dependencies, including
conditions, negations, and temporal qualifiers. Its output records which checked
items are retained, contradicted, or unresolved and links them to supporting
source spans. CertMem then selects a compact \emph{evidence residual}: verbatim
passages intended to repair the most consequential omissions. Selection uses
predicted reduction in answering loss per added token and reassesses overlapping
candidates as the residual grows. When a reference or condition requires several
passages jointly, the candidate is a passage bundle. The resulting memory
contains the summary, selected residual, source links, and a bounded audit record;
the audit describes inspected content rather than certifying exhaustive coverage.

At query time, a retriever first supplies compact memories and their residuals.
The controller relates the query to the recorded gaps and selects additional
source spans when their predicted benefit justifies the added cost. Each expansion
is followed by reassessment, so reading part of a source does not automatically
mark all its omissions as resolved. A bounded search over the raw archive remains
available when relevant evidence lies outside the retrieved summary links.
The controller stops when evidence passes a validation-selected threshold or
the read budget is exhausted, recording unresolved cases. Development uses paired answers from
compact and expanded contexts with a fixed reader. Future benchmark questions
are excluded from memory construction, and policy selection uses conversations
separate from the final evaluation. Neither an entailment score nor a stopping
decision is treated as a guarantee that an individual answer is correct.

Our optimization target makes the full cost of this procedure explicit. For a
history $H$, queries $q_{1:N}$, and a write-and-read policy $\pi$, we consider
\begin{equation}
\min_{\pi}\;
\mathbb{E}\!\left[
 C_{\mathrm{maint}}^{\pi}(H)
 + \sum_{j=1}^{N} C_{\mathrm{query}}^{\pi}(H,q_j)
\right]
\quad\text{subject to}\quad
R_{+}(\pi)\leq\epsilon .
\label{eq:objective}
\end{equation}
Here, maintenance includes compression, auditing, and residual construction;
query cost includes routing, evidence inspection, and answer generation.
$R_{+}$ is the expected positive increase in task loss over a fixed source-based
reference, using the same ground truth. Absolute accuracy is also required because
the reference can fail. Equation~\ref{eq:objective} specifies a target rather than
an established guarantee. Persistent storage, including the raw archive and
metadata, is reported separately.

We outline evaluation on LoCoMo and LongMemEval under matched reader models,
data splits, and token budgets. The decisive comparisons are with summary-only
memory, raw retrieval, the five preceding frameworks, and controlled variants
that remove residuals or replace their selection with relevance or fact-coverage
scores. Accuracy--cost curves, false stopping decisions, and repeated-query
workloads test whether proactive retention offsets its maintenance overhead.
The proposed contributions are: (i) a formulation linking write-time retention
and query-time evidence allocation through recorded compression omissions;
(ii) an evidence-residual mechanism that uses these omissions for both selective
preservation and source expansion; and (iii) an evaluation protocol that tests
this mechanism through downstream answer quality and complete token accounting.
Empirical improvements remain to be measured.

\bibliography{references}
\bibliographystyle{iclr2027_conference}
\end{document}
